%% 
\documentclass[preprint,1p]{article}
%\usepackage[left=1cm, right=1cm, top=2cm, bottom=2cm]{geometry}
\usepackage{geometry}
\usepackage{hyperref}

%% Use the option review to obtain double line spacing
%% \documentclass[authoryear,preprint,review,12pt]{elsarticle}

%% Use the options 1p,twocolumn; 3p; 3p,twocolumn; 5p; or 5p,twocolumn
%% for a journal layout:
%% \documentclass[final,1p,times]{elsarticle}
%% \documentclass[final,1p,times,twocolumn]{elsarticle}
%% \documentclass[final,3p,times]{elsarticle}
%% \documentclass[final,3p,times,twocolumn]{elsarticle}
%% \documentclass[final,5p,times]{elsarticle}
%% \documentclass[final,5p,times,twocolumn]{elsarticle}

%% For including figures, graphicx.sty has been loaded in
%% elsarticle.cls. If you prefer to use the old commands
%% please give \usepackage{epsfig}

%% The amssymb package provides various useful mathematical symbols
\usepackage{amssymb}
\usepackage{epsfig}
%% The amsthm package provides extended theorem environments
%% \usepackage{amsthm}

%% The lineno packages adds line numbers. Start line numbering with
%% \begin{linenumbers}, end it with \end{linenumbers}. Or switch it on
%% for the whole article with \linenumbers.
%% \usepackage{lineno}
%\usepackage{CJK}
\usepackage{amssymb}
\usepackage{latexsym,bm}
\usepackage[tbtags]{amsmath}
\usepackage{amssymb}
\usepackage{graphicx,subfigure}
\usepackage{enumerate} 
\usepackage{fancyhdr}
\usepackage{float}
\usepackage{cite}
\usepackage{tabularx}
\usepackage[flushleft]{threeparttable}
\usepackage{booktabs}
\usepackage{dcolumn}
\usepackage{graphicx}
\usepackage{wrapfig}%(P227)
\usepackage{multicol}
\usepackage{verbatim}
\usepackage{mathrsfs}

\usepackage{multirow}
\usepackage{IEEEtrantools}
\usepackage{graphicx}


\usepackage{listings}
\usepackage{url}
\usepackage{amsmath}
\usepackage{amsfonts}
\usepackage{amssymb}
\usepackage{graphicx}
\usepackage{algorithm}
\usepackage{algpseudocode}
\usepackage{setspace}
\usepackage{listings}
\usepackage{array}
\usepackage[dvipsnames]{xcolor}
\usepackage{todonotes}
\newcommand{\TODO}[2][]{\todo[inline, #1]{#2}}

\renewcommand{\sfdefault}{phv}                         
\renewcommand{\rmdefault}{ptm}
\renewcommand{\ttdefault}{pcr}
\renewcommand{\normalfont}{}



\newcommand{\place}[2]{ 
	\vspace{0.5cm}
	\begin{center}
		{\bf \large Place  #1 \ref{#2} here }
	\end{center}
	\vspace{0.5cm}
}


\newcommand{\makenumber}[2]{ 
	\newpage
	%\thispagestyle{empty}
	%\adDtocounter{page}{-1}
	%{\bf \large Author: {\it Jose Miguel Mantas Ruiz}}
	%{\bf \large #1 \ref{#2}}
	%\newpage
}


\newcommand{\tablecaption}[1]{\centering{\bf #1} \vspace{-0.5cm} }



\usepackage{float}
\newfloat{algorithm}{t}{lop}
\DeclareMathOperator{\arcsinh}{arcsinh}
\DeclareMathOperator{\arccosh}{arccosh}
\DeclareMathOperator{\arctanh}{arctanh}
\newcolumntype{L}[1]{>{\raggedright\let\newline\\\arraybackslash\hspace{0pt}}m{#1}}
\newcolumntype{C}[1]{>{\centering\let\newline\\\arraybackslash\hspace{0pt}}m{#1}}
\newcolumntype{R}[1]{>{\raggedleft\let\newline\\\arraybackslash\hspace{0pt}}m{#1}}
\newcommand{\pythode}{\texttt{pythODE}}
\newcommand{\pypp}{\texttt{pythODE++}}
\newcommand{\rodas}{RODAS}
\newcommand{\radau}{RADAU5}
\newcommand{\adolc}{\texttt{ADOL-C}}
\newcommand{\umfpack}{\texttt{UMFPACK}}
\newcommand{\colpack}{\texttt{ColPack}}
\newcommand{\cusp}{\texttt{CUSP}}
\newcommand{\order}[1]{\mathcal O \big( #1 \big)}
\newcommand{\by}{\mathbf y}
\newcommand{\byt}{\mathbf{\tilde y}}
\newcommand{\bphi}{\boldsymbol \Phi}
\newcommand{\bpsi}{\boldsymbol \Psi}
\newcommand{\deriv}[2]{\frac{\textrm{d} #1}{\textrm{d} #2}}
\newcommand{\derivn}[3]{\frac{\textrm{d}^{#3} #1}{\textrm{d} #2^{#3}}}
\newcommand{\pderiv}[2]{\frac{\partial #1}{\partial #2}}
\newcommand{\pderivn}[3]{\frac{\partial^{#3} #1}{\partial #2^{#3}}}
\newcommand{\mpd}[3]{\frac{\partial^2 #1}{\partial #2 \partial #3}}
\newcommand{\Dt}{\Delta t}

\newcommand{\bs}{\textbackslash{}}



\newcommand{\NEW}[1]{\textcolor{black}{#1}}
\newcommand{\RevOne}[1]{\textcolor{black}{#1}}
\newcommand{\NOTE}[1]{\fcolorbox{yellow}{red}{#1}}

%\newenvironment{NEWB}  {\begin{color}{red}} { \end{color}}

\newcommand{\RevTwo}[1]{\textcolor{black}{#1}}

\newenvironment{NEWB}{\begin{color}{black}} { \end{color}}

\begin{document}

%

\title{ \bf User Manual for the VSIIE solvers}

\maketitle 

\begin{abstract}

This document provides a user manual for a software framework 
that implements different Variable Stepsize Implicit-Implicit-Explicit (VSIIE) solvers.
The Implicit-Implicit-Explicit (IIE) solvers \cite{MARABEH2023} are a family of high-order time integration methods designed
to efficiently and accurately solve stiff systems of ordinary differential equations (ODEs) that often
arise from the spatial discretization of relevant models based on partial differential equations (PDEs)
such as diffusion-reaction-advection (DRA) models, which can be divided into three separate terms.
Here, the IIE solvers are extended to allow for variable stepsize, resulting in the VSIIE solvers.
In addition, an implementation of Variable Stepsize Semi-implicit Backward Differentiation Formula (VSSBDF) solvers 
 \cite{MARABEH2025} is also included for the purposes of experimental comparison.

\end{abstract}


\section{User manual}

The user manual describes how to use the software framework and several \texttt{Python} scripts 
developed to perform automatically experiments to compare the performance of the 
fixed stepsize IIE  methods \cite{MARABEH2023},the  variable stepsize VSSBDF methods \cite{MARABEH2025}, and the variable stepsize 
VSIIE  methods \cite{MARABEH2026}.

\subsection{Building and Calling the solver}

The executable of the solver can be generated by typing {\tt make} in the main folder (usually named {\tt VSIIE}) 
where the source code is located.  This will generate the executable file called {\tt IVP\_Solver} in the same folder.

The solver can be called from the command line in a terminal window.  The call follows the following format:

\vspace{0.3cm}
\framebox[1.1\width]{\centering

 {\bf IVP{\_}Solver} 
 {\em Solver \hfill problem{\_}id. \hfill order \hfill Neqn \hfill tf \hfill stepsize \hfill tol \hfill num{\_}reps \hfill
  alpha \hfill eta{\_}min  \hfill eta{\_}max}	
}
\vspace{0.3cm}

\noindent where:

\begin{itemize}

    \item {\em Solver}: It denotes the type of solver to be used. 
        The provided values can be:
	    \begin{itemize}
			\item {\tt IIE}: A constant stepsize IIE method is used.
            \item {\tt VSIIE}: A Variable Stepsize IIE (VSIIE) method is used.
            \item {\tt VSSBDF}: A Variable Stepsize Semiimplicit BDF (VSSBDF) method is used.
	    \end{itemize} 

	\item {\em problem\_id.}: It denotes the mathematical model to be solved. The provided values can be:
	    \begin{itemize}
		    \item {\tt 0}: 1D Advection-Diffusion.	
		    \item {\tt 1}: Stiff Brusselator.
		    \item {\tt 2}: FKPP Combustion.		
		    \item {\tt 3}: Ignition Combustion.
		    \item {\tt 4}: Fisher Combustion.
		    \item {\tt 5}: High Irradiance RESponse (HIRES) problem \cite{HIRES}.
		\end{itemize}

    \item {\em order}: It denotes the order of the VSIIE/IIE method to be used (1, 2, 3  or 4).

	\item {\em Neqn}: It denotes the number of ODEs in the system to be solved 
	(It is important not to confuse this with the number of grid points.).
	
	\item {\em tf}: It denotes the final simulation time. 

    \item {\em stepsize}: It denotes the constant stepsize if const\_var=const (IIE method)  
	                 and the initial stepsize if const\_var=var (VSIIE method)

    \item {\em tol}: Error tolerance which it is only used for VSIIE (const\_var=var)

    \item {\em num\_reps}:  Number of repetitions of the experiment (the experiment is repeated several times in oprder 
	to obtain a report the minimum runtime of all repetitions).

    \item {\em alpha eta\_min eta\_max} (Optional):  Parameters $\alpha$, $\eta_{\min}$ and $\eta_{\max}$ of 
	the adaptive time-stepping strategy (by default: $\alpha=0.8$, $\eta_{\min}=0.5$ and $\eta_{\max}=4$).

\end{itemize}

Therefore, each call to the solver is applied to solve a particular IVP model instance using a particular method 
(VSIIE, VSSBDF or IIE) of a given order.

When the solver is called, generates output by the screen and in output files.
The first time the solver is called for a particular problem and final time, 
it creates a folder called {\tt data} in the same folder where the executable is located and
generates a text file where the reference solution for this problem instance is saved. 
This reference solution is computed by using the VSSBDF method of order 4 with a very small tolerance of 
$10^{-6}$. Thus, the reference solution
data can be reused in subsequent calls of the solver for the same problem instance, avoiding recomputation.

After running the solver, it generates a line with data in a text output file with {\tt .txt} extension, whose name
depends on the particular IVP model, method and final time. This output file is located in the same folder of the executable.
The output file contains the following data in each line (separated by several spaces):
\begin{itemize}
\item {\bf For VSIIE and VSSBDF}: [tolerance] \quad [error] \quad  [Minimum runtime] \quad  [Number of time steps] 
\quad  [Number of rejected time steps] \quad [Number of accepted time steps]
\quad [Average number of Newton iterations per implicit solve] 

\item {\bf For IIE}:  [Constant stepsize] \quad [Error] \quad [Minimum runtime] \quad [Number of time steps]
\quad  [Number of rejected time steps=0] \quad [Number of accepted time steps]
\quad [Average number of Newton iterations per implicit solve] 

\end{itemize}


\subsection	{The scripts}	
There is a \texttt{Python} script provided to run automatically the experiments corresponding to each particular mathematical IVP model 
and generate the corresponding plots (work-precision and time-steps).
Each script runs the experiments for the IIE, VSSBDF and VSIIE  methods using different parameters. Each \texttt{Python} script must be called 
in the particular folder where the script is located (although the user can call it from other locations if he creates a new folder, copies the script 
to the new folder and specifies what is the name of the new folder in the script). The location and correspondin model of the particular scripts 
are given in Table \ref{Scripts_table}.



\begin{table}[h!]
	\caption{Data about the particular python scripts}{\label{Scripts_table}}
	\centering 
			\begin{tabular}{|c|c|c|}
			\hline
			{\bf Script Name}       & {\bf Folder}       & {\bf Model   } \\ \hline
			{\tt adr.py}      & {\tt ADR}    & 1D Advection-Diffusion \\ \hline
			{\tt sbr.py}      & {\tt SBR}    & Stiff Brusselator \\ \hline
			{\tt fkpp.py}     & {\tt FKPP}   & FKPP Combustion \\ \hline
			{\tt ignition.py} & {\tt IGNIT}  & Ignition Combustion \\ \hline
			{\tt fisher.py}   & {\tt FISHER} & Fisher Combustion \\ \hline
			{\tt hires.py}   & {\tt HIRES}   & High Irradiance RESponse (HIRES) problem \\ \hline
		\end{tabular}
\end{table}


For each \texttt{Python} script, the user must define the following variables 
which define the experiments to be run:

\begin{itemize}
\item {\tt Problem}: It denotes the problem id (0 to 4) corresponding to the mathematical IVP model.
\item {\tt IVP\_name}: It denotes the name of the mathematical IVP model. The provided values can be:
	    \begin{itemize}
		    \item {\tt "1D\_Adv-Diff"}: 1D Advection-Diffusion.	
		    \item {\tt "stiff\_Brusselator"}: Stiff Brusselator.
		    \item {\tt "FKPP"}: FKPP Combustion.		
		    \item {\tt "Ignition"}: Ignition Combustion.
		    \item {\tt "Fisher"}: Fisher Combustion.
		    \item {\tt HIRES}: High Irradiance RESponse (HIRES) problem.
		\end{itemize}
\item {\tt Neqn}: It denotes the number of ODEs in the system to be solved 
  (It must be the the number of equations in the chosen IVP model, not the number of grid points).

\item {\tt Tf}: It denotes the final simulation time.

\item {\tt FOLDER}: It denotes the folder name where the experiments will be run 
  (It must be the folder where the chosen \texttt{Python} script is located). For instance, "ADR" for the 1D Advection-Diffusion model.

\item {\tt stepsize\_array}: It denotes the 2D array of constant stepsizes to be used for the IIE methods of different orders. 
It is also used to define the initial stepsizes for the VSIIE methods of different orders.
For instance, {\tt stepsize\_array[0]} contains the array of constant stepsizes for the fisrst order IIE method and the initial 
stepsizes for the first order VSIIE method.

\item {\tt VSIIE\_tol\_array}: It denotes the array of error tolerances to be used for the VSIIE methods of different orders.
For ADR, the definition of this array can be as follows: 



\begin{small}
$stepsize\_array=[$
\begin{align*}
& [7.5e-3, 1e-3, 5e-4, 1.0e-4, 5e-5, 1e-5, 5e-6, 1e-6, 5e-7, 1e-7], \#order 1\\
& [5e-2, 2.5e-2, 1e-2, 7.5e-3, 5e-3, 1e-3, 5e-4, 1e-4, 5e-5, 1e-5, 5e-6, 1e-6],  \#order 2\\
& [1e-1, 2.5e-2, 1e-2, 7.5e-3, 5e-3, 1e-3, 5e-4, 1e-4],  \#order 3\\
& [1e-1, 7.5e-2, 5e-2, 2.5e-2, 1e-2, 7.5e-3, 5e-3, 1e-3, 5.0e-4]]   \#order 4	
\end{align*}
\end{small}


\item{\tt VSSBDF\_tol\_array}: It denotes the array of error tolerances to be used for the VSSBDF 
methods of different orders. It is defined in the same way as {\tt VSIIE\_tol\_array}. 


\item {\tt alpha\_array}: It denotes the array of values which indicates the $\alpha$ parameter values (which affects to the the time-stepping 
algorithm) to be used for each particular experiment of the  VSIIE methods of different orders.
For ADR, the definition of this array can be as follows:

\begin{small}
$alpha\_array=[$
\begin{align*}
	&[0.8, 0.8, 0.8, 0.8, 0.8, 0.8, 0.8, 0.8, 0.8, 0.8],   \#order 1\\
	&[0.4, 0.4, 0.4, 0.6, 0.6, 0.6, 0.6, 0.6, 0.6, 0.6, 0.6, 0.6], \#order 2\\
	&[0.3, 0.3, 0.4, 0.4, 0.4, 0.4, 0.4, 0.4], \#order 3 \\
	&[0.3, 0.3, 0.3, 0.3, 0.3, 0.3, 0.3, 0.3, 0.3 ]]  \#order 4			
\end{align*}
\end{small}

\end{itemize}

The arrays {\tt stepsize\_array}, {\tt VSIIE\_tol\_array} and {\tt alpha\_array} must have the same number of elements for each order.


\subsection{Calling the \texttt{Python} scripts}
Each \texttt{Python} script can be called from the command line in a terminal window (being on the folder where the script is located).  
The call has the following format:


\vspace{0.3cm}
\framebox[1.1\width]{\centering
{\bf script\_name.py}  \hfill  {\em run\_type} \hfill  {\em min\_order} \hfill {\em max\_order} \hfill {\em IIE\_min\_order} \hfill {\em IIE\_max\_order} \hfill {\em N\_Reps}	
}

\vspace{0.3cm}
\noindent where:

\begin{itemize}

\item {\em run\_type}: It denotes if the experiments are to be run or not. The provided values can be:
	    \begin{itemize}
			\item {\tt 0}: Do not run the experiments but only plot them from the previously obtained data files.
			\item {\tt 1}: Run the experiments, save the results to data files and then plot them.
	    \end{itemize}
\item	{\em min\_order}   {\em max\_order}: Minimum and maximum order of the corresponding VSIIE/VSSBDF family of methods 
        to be used in the experiments. 


\item	{\em IIE\_min\_order}   {\em IIE\_max\_order} : Minimum and maximum order of the IIE methods to be used in the experiments.

\item {\em N\_Reps}: Number of repetitions of each experiment to obtain the minimum runtime.	

\end{itemize}

If the user	wish to run the experiments for all the orders of all VSIIE, VSSBDF and IIE methods with 3 repetitions,
he can call the script as follows: 

\begin{center}
{\bf script\_name.py} \hspace{2mm} 1\hspace{2mm} 1 \hspace{2mm} 4 \hspace{2mm} 1 \hspace{2mm} 4 \hspace{2mm} 3.
\end{center}

On the other hand, if the user wish to run the experiments for a particular range of orders of VSIIE and/or IIE methods,
he can specify the corresponding minimum and maximum orders in the call to the script.
So, for example, the following call


\begin{center}
{\bf script\_name.py} \hspace{2mm} 1 \hspace{2mm} 2 \hspace{2mm} 3 \hspace{2mm} 1 \hspace{2mm} 2 \hspace{2mm} 4	
\end{center}


\noindent runs the experiments for orders 2 and 3 of the VSIIE/VSSBDF methods and orders 1 and 2 of the IIE methods using 4 repetitions.

One can also avoid easily to do experiments with a particular method (IIE or VSIIE). So for example, the 
the call:


\begin{center}
{\bf script\_name.py} \hspace{2mm}  1 \hspace{2mm} 2 \hspace{2mm} 2 \hspace{2mm} 0 \hspace{2mm} 0 \hspace{2mm} 3 
\end{center}

\noindent only runs the experiments for order 2 of the VSIIE methods (there is no experiment for IIE methods). 
In this way, the scripts allows the user to choose what orders to be used in the experiments and thus avoid to recompute 
data which have been already obtained previously.

\bibliographystyle{plain}
\bibliography{references}

\end{document}
